\chapter{Brain Disorders}

\section*{Definition and Overview}
The term \emph{brain disorders} describes the broad range of structural, functional, developmental, and psychological conditions that affect the brain and its workings. It overlaps with, but is wider than, the concept of \emph{brain diseases}: as well as stroke, dementia, infection, and tumour, it embraces neurodevelopmental conditions such as autism spectrum disorder and attention deficit hyperactivity disorder (ADHD), epilepsy, migraine and other headache disorders, movement disorders, the consequences of brain injury, and mental health conditions that arise from altered brain function, such as major depression, bipolar disorder, and psychosis. What these disparate conditions share is disruption of the brain's normal wiring, chemistry, or electrical activity. No single list is exhaustive, and some conditions -- notably autism and ADHD -- are increasingly understood as neurodevelopmental differences (variations in functioning) rather than diseases in the conventional sense, a shift with important implications for how they are assessed, supported, and destigmatised.

\section*{Epidemiology}
Brain disorders, considered as a group, are among the commonest health problems in the world. Headache disorders affect a large proportion of the population, with migraine alone estimated to affect around 1 in 7 people. Epilepsy has a lifetime prevalence of roughly 1 in 100, with onset most often in childhood or after the age of 60. ADHD affects around 5--8\% of children and often persists into adulthood, while autism is now diagnosed in roughly 1--2\% of children in many countries, reflecting improved recognition as much as any true change in frequency. Dementia affects about 5--8\% of people over 65 and rises steeply with age. Mental health disorders are among the leading causes of disability worldwide, and because the category is so broad, most people will have personal or family contact with a brain disorder during their lifetime.

\section*{Aetiology and Pathophysiology}
The causes of brain disorders are diverse and are often multiple and interacting:

\begin{itemize}
    \item \textbf{Genetic and developmental:} variation in genes and in early brain development underlies autism, ADHD, many epilepsy syndromes, and some forms of intellectual disability
    \item \textbf{Structural:} tumours, congenital malformations, and the consequences of traumatic injury
    \item \textbf{Neurochemical:} altered signalling between neurones, mediated by neurotransmitters such as serotonin, dopamine, and glutamate, is implicated in mood, anxiety, and psychotic disorders
    \item \textbf{Electrical:} abnormal, synchronised neuronal discharge produces the seizures of epilepsy
    \item \textbf{Vascular and metabolic:} reduced blood flow, hypoxia, hypoglycaemia, and toxins disturb cerebral function
    \item \textbf{Degenerative:} progressive neuronal loss underlies the dementias and Parkinson's disease
    \item \textbf{Infectious and inflammatory:} infections and immune-mediated attacks, including autoimmune encephalitis, damage brain tissue
\end{itemize}

In many cases several factors combine -- for example, a genetic predisposition plus environmental stressors may together trigger a mood disorder or a first seizure.

\section*{Clinical Features}
The presentation depends on the disorder and the brain regions involved, but common patterns include:

\begin{itemize}
    \item \textbf{Cognitive difficulties:} problems with attention, memory, planning, language, or learning
    \item \textbf{Seizures:} generalised convulsions, focal attacks, or staring spells and unexplained blackouts
    \item \textbf{Headache:} including migraine with or without aura and other recurrent headache syndromes
    \item \textbf{Mood disturbance:} persistent low mood, anxiety, irritability, or episodes of elevated mood
    \item \textbf{Perceptual and thought disturbance:} hallucinations, delusions, or disorganised thinking
    \item \textbf{Movement problems:} tremor, tics, clumsiness, slowness, or involuntary movements
    \item \textbf{Behavioural change:} impulsivity, difficulty with social interaction, or altered personality
\end{itemize}

\section*{Differential Diagnosis}
The overlap between disorders means that careful assessment is essential:

\begin{itemize}
    \item Structural disorders (tumour, stroke, injury) presenting with psychiatric symptoms versus primary mental health conditions
    \item Epilepsy versus syncope, migraine, or psychogenic non-epileptic seizures
    \item Depression versus dementia versus hypothyroidism or vitamin B12 deficiency in older people
    \item ADHD versus anxiety, sleep deprivation, or the effects of substance use
    \item Parkinsonism versus drug-induced tremor or functional neurological disorder
    \item Delirium (an acute, reversible confusion from a physical cause) versus dementia (a gradual decline)
    \item Autism versus social anxiety or language disorder in children
\end{itemize}

\section*{When to Seek Medical Care}
See a doctor for new or worsening problems with memory, mood, seizures, headaches, movement, or behaviour, particularly when they interfere with work, school, or relationships. Children with developmental or behavioural concerns benefit from early assessment, which opens the door to support and intervention. A sudden change in mental state or behaviour always warrants prompt review, as physical causes must be excluded.

\textbf{Call 000 (or local emergency services) for:}
\begin{itemize}
    \item A seizure lasting more than five minutes, or repeated seizures without recovery between them
    \item Sudden weakness, numbness, difficulty speaking, or confusion -- possible stroke
    \item Loss of consciousness or a sudden, severe headache
    \item Deliberate self-harm or talk of suicide
    \item Severe agitation, hallucinations, or behaviour that puts the person or others at risk
\end{itemize}

\section*{Diagnosis and Investigations}
\begin{itemize}
    \item \textbf{History and examination:} detailed accounts from the person and, where relevant, family members or teachers, supplemented by a full physical and neurological examination
    \item \textbf{Structured interviews and rating scales:} standardised tools that assess mood, anxiety, attention, and development, and aid diagnosis and monitoring
    \item \textbf{Blood tests:} thyroid function, glucose, vitamin levels, inflammatory markers, and other checks to exclude physical causes
    \item \textbf{Imaging:} CT or MRI where a structural, vascular, or inflammatory cause is suspected
    \item \textbf{Electroencephalography (EEG):} records cerebral electrical activity and is central to diagnosing and classifying epilepsy
    \item \textbf{Neuropsychological assessment:} formal testing of memory, attention, and executive function
    \item \textbf{Psychiatric assessment:} specialist evaluation where a mental health condition is possible, and joint neurology--psychiatry review in complex cases
\end{itemize}

\section*{Management}
Most brain disorders are managed long term, typically by a team that may include a neurologist, psychiatrist, psychologist, and allied health professionals, working with the person and their family:

\begin{itemize}
    \item \textbf{Medicines:} antiepileptic drugs, antidepressants, mood stabilisers, antipsychotics, stimulants and non-stimulants for ADHD, and agents for movement disorders
    \item \textbf{Psychological therapies:} cognitive behavioural therapy and other talking treatments for mood, anxiety, and many functional disorders
    \item \textbf{Behavioural and educational support:} structured routines, classroom accommodations, and evidence-based behavioural programmes for ADHD and autism
    \item \textbf{Neuromodulation:} transcranial magnetic stimulation, vagus nerve stimulation, or deep brain stimulation for selected treatment-resistant conditions
    \item \textbf{Rehabilitation:} physiotherapy, occupational therapy, and speech and language therapy to address specific impairments
    \item \textbf{Psychosocial support:} family education, peer support groups, advocacy, and practical help with daily living and employment
\end{itemize}

\section*{Complications}
\begin{itemize}
    \item Injuries and accidents related to seizures, impaired attention, or poor judgement
    \item Difficulty at school, at work, and in relationships
    \item Substance misuse as a way of coping with symptoms
    \item Social isolation and stigma
    \item Worsening of physical health through inactivity, poor sleep, or neglected self-care
    \item Increased risk of self-harm and suicide, particularly in mood and psychotic disorders
    \item Adverse effects of medication and financial or carer strain on families
\end{itemize}

\section*{Prognosis and Outcomes}
Outcomes vary enormously. Many disorders -- including most epilepsies, mood and anxiety disorders, and ADHD -- respond well to treatment, and most affected people live full, productive lives. Neurodevelopmental conditions such as autism are lifelong, but with appropriate support many individuals flourish, and diagnosis can be a source of identity and empowerment. Degenerative disorders such as dementia progress over time, though the pace varies greatly and treatment can slow decline and preserve quality of life. Across all brain disorders, early diagnosis, evidence-based treatment, and strong social support are consistently associated with better outcomes, and remission or substantial improvement is realistic for the majority of treatable conditions.

\section*{Prevention and Self-Care}
\begin{itemize}
    \item Protect the brain from injury with seatbelts and helmets for cycling, riding, and contact sport
    \item Prioritise good sleep, regular physical activity, and a balanced diet
    \item Manage stress and seek help early for mood, anxiety, or behavioural problems
    \item Avoid heavy alcohol and recreational drug use, especially if you have a known seizure tendency
    \item Keep up to date with developmental and behavioural assessments in children
    \item Take prescribed medicines reliably and attend scheduled reviews for chronic conditions
    \item Build and maintain social connections, and draw on family, friends, and support organisations
\end{itemize}

\section*{Sources and Further Reading}
The following organisations provide reliable, evidence-based information for patients, students, and clinicians.

\begin{itemize}
    \item NHS. \textit{Brain and nervous system disorders: patient information.} https://www.nhs.uk/conditions/
    \item Mayo Clinic. \textit{Brain, mental health, and neurodevelopmental conditions: patient education.} https://www.mayoclinic.org/diseases-conditions/
    \item World Health Organization. \textit{Mental and neurological disorders: global information and policy.} https://www.who.int/
    \item Centers for Disease Control and Prevention. \textit{Brain health, epilepsy, and developmental conditions: public health resources.} https://www.cdc.gov/
    \item MSD Manual. \textit{Neurologic disorders and mental health disorders: consumer and professional reference.} https://www.msdmanuals.com/
    \item MedlinePlus. \textit{Brain diseases and neurological conditions: consumer health information.} https://medlineplus.gov/
\end{itemize}
